\documentclass[conference]{IEEEtran}
\IEEEoverridecommandlockouts
\usepackage{cite}
\usepackage{amsmath,amssymb,amsfonts}
\usepackage{graphicx}
\usepackage{textcomp}
\usepackage{xcolor}
\usepackage{booktabs}
\usepackage{multirow}
\usepackage{url}
\usepackage{hyperref}

\def\BibTeX{{\rm B\kern-.05em{\sc i\kern-.025em b}\kern-.08em
    T\kern-.1667em\lower.7ex\hbox{E}\kern-.125emX}}

\begin{document}

\title{Climate-Smart Crop Calendar Recommendation and Monsoon Shift\\
Analysis for Tamil Nadu Using Machine Learning}

% TODO: replace with your actual name(s), department, institution, city, country
\author{
\IEEEauthorblockN{Author Name}
\IEEEauthorblockA{\textit{Department of Computer Science} \\
\textit{Your College/University Name} \\
City, India \\
your.email@example.com}
}

\maketitle

\begin{abstract}
Most machine-learning crop recommendation systems in the literature are trained and
evaluated on small, generic, non-regional datasets (soil N-P-K, temperature, humidity,
pH, and rainfall) and treat rainfall as a single static feature rather than analyzing
how a region's monsoon behavior is actually shifting over time; they are also rarely
carried through to a deployed, farmer-facing system. This paper presents a
climate-smart crop calendar system for Tamil Nadu, India, built around two tightly
coupled components: (i) district-level monsoon shift analysis, using each district's
real 21-year historical rainfall record to detect the direction and magnitude of
long-term rainfall change and translate it into an explicit drought/flood risk
signal, and (ii) a machine-learning yield and sowing-window recommendation pipeline
that is conditioned on that district's climate history rather than a generic
national average. The system is built from real district-level government
agricultural census data (38 districts, 25 crops, six government-defined cropping
seasons, historical records spanning 2005--2019) merged with district-level
historical climate data, combined with a TNAU-style agronomic sowing-window reference
calendar. While preparing this system we discovered a systematic $\sim$100$\times$
production-unit inconsistency present in the underlying government yield records for
every crop in the dataset -- silently corrupting a substantial fraction of the
training signal used by both the yield model and the rainfall-trend analysis -- and
we describe an empirical, duplicate-record-grounded procedure to detect and correct
it (6{,}295 of the affected records were rescaled; one record with an unrecoverable
Area value was dropped). After this correction, we trained and fairly compared three
regression models -- a decision tree, a random forest, and XGBoost -- for
district-level yield estimation, with hyperparameters tuned under matched budgets
(tree count $\times$ depth) rather than leaving the gradient-boosted model
under-regularized relative to the bagged baseline. The tuned XGBoost model achieved
the best held-out performance (test $R^2 = 96.98\%$, MAE $= 1.45$, RMSE $= 3.82$)
against Random Forest ($R^2 = 96.57\%$) and a single Decision Tree ($R^2 = 95.27\%$),
confirmed with 5-fold cross-validation. The resulting recommendation pipeline --
climate retrieval, monsoon-shift-aware drought/flood risk assessment, yield
prediction, sowing-window selection, and a data-driven confidence score -- is
deployed as a full-stack web application (Flask, PostgreSQL) that is also installable
as a Progressive Web App and as a native Android application, with real-time weather
integration replacing static climate estimates for the farmer-facing dashboard. We
discuss the practical significance of the data-quality finding, the deployed system's
architecture, and outline uncertainty quantification as a direction for future work.
\end{abstract}

\begin{IEEEkeywords}
crop recommendation, monsoon shift analysis, rainfall trend, precision agriculture,
machine learning, XGBoost, data quality, yield prediction, deployed agricultural
decision support system
\end{IEEEkeywords}

\section{Introduction}

Machine learning (ML) has become a standard tool for crop recommendation and yield
prediction, supporting farmers with data-driven guidance on what to grow and when
\cite{alam2025}. A large share of published crop recommendation studies, however, train
and validate models on a single widely-used public dataset of 2{,}200 samples spanning
22 crop classes with seven generic agro-climatic features (nitrogen, phosphorus,
potassium, temperature, humidity, pH, and rainfall) \cite{alam2025}. This dataset is
convenient and well-balanced, but it is not tied to any specific region's actual
administrative and agronomic structure -- it has no district, no government-defined
cropping season, and no historical yield trend for a real place a farmer actually
farms in. Consequently, models trained on it, however accurate on their own test split,
cannot directly tell a farmer in a named district when to sow a named crop this season.

A second, related gap is that even studies which do use real rainfall data typically
treat it as one more static input feature at prediction time, rather than first asking
whether the region's monsoon behavior itself has been \emph{shifting} over the
available historical record. Trend-detection studies of Indian rainfall (e.g.,
\cite{praveen2020}), by contrast, show that this question has a real, non-trivial
answer -- most Indian meteorological divisions exhibit statistically significant
negative rainfall trends over the twentieth century -- and that answer is exactly the
kind of district-specific signal a crop calendar recommendation should be conditioned
on, not just a raw rainfall number for the current season.

This paper takes both gaps as its starting point: instead of a generic feature-based
dataset and a static rainfall feature, we build a crop recommendation and
sowing-calendar system directly from government agricultural census and
district-level climate records for the Indian state of Tamil Nadu, explicitly analyze
each district's real multi-decade rainfall trend as a monsoon-shift signal that feeds
the recommendation's risk assessment, and carry the resulting model through to a
deployed, installable application. This choice surfaced a data problem that a small
curated benchmark dataset would not: real government production records, aggregated
over decades from multiple reporting cycles, contain unresolved unit inconsistencies.
We detail this issue and our correction procedure as one of this paper's
contributions, since it directly affects the validity of any yield model \emph{and}
any rainfall-trend analysis trained on this class of data source, and, to our
knowledge, is not discussed in prior crop recommendation literature.

The contributions of this paper are:
\begin{enumerate}
  \item A district-level monsoon shift analysis over each Tamil Nadu district's real
        21-year historical rainfall record, translating the detected trend direction
        and magnitude into an explicit drought/flood risk signal that conditions the
        downstream crop recommendation, rather than treating rainfall as a single
        static feature.
  \item A district-level crop sowing recommendation pipeline built on real Tamil Nadu
        government agriculture and climate records (25 crops, 38 districts, six
        government-defined seasons), combining a TNAU-style agronomic reference
        calendar with the monsoon-shift-based risk assessment above.
  \item Identification and correction of a systematic $\sim$100$\times$ production-unit
        inconsistency affecting every crop in the underlying government dataset, using
        an empirical procedure grounded in confirmed duplicate-record pairs rather than
        a purely statistical outlier heuristic.
  \item A fair, matched-budget hyperparameter comparison of Decision Tree, Random
        Forest, and XGBoost regressors for yield estimation on the corrected data,
        validated with both a held-out split and 5-fold cross-validation.
  \item A deployed, farmer-facing full-stack implementation (Flask, PostgreSQL) that is
        installable as a Progressive Web App and as a native Android application, with
        live weather data replacing static estimates on the dashboard.
\end{enumerate}

The remainder of this paper is organized as follows. Section II reviews related work
in ML-based crop recommendation and rainfall/monsoon trend analysis. Section III
describes the data source, the crop inclusion criterion, and the production-unit
correction procedure. Section IV describes the monsoon shift analysis and the
recommendation pipeline and model configurations. Section V presents the model
comparison results. Section VI describes the deployed system. Section VII discusses
limitations and future work, followed by the conclusion.

\section{Related Work}

Recent crop recommendation research has explored ensemble learning, IoT integration,
and cloud deployment to improve prediction reliability and reach. Alam et
al.~\cite{alam2025} proposed a soft-voting ensemble of Logistic Regression, Na\"ive
Bayes, KNN, Decision Tree, Random Forest, and XGBoost on the 2{,}200-sample generic
crop dataset described above, reporting 99.54\% classification accuracy and
introducing entropy-based uncertainty quantification over the ensemble's softmax
probabilities -- to our knowledge the first crop recommendation study to explicitly
quantify prediction confidence rather than deliver a purely deterministic label.
Hasan et al.~\cite{hasan2023} similarly used a K-nearest-neighbor/Random
Forest/Ridge-regression ensemble for crop production prediction in Bangladesh, and
Bandi et al.~\cite{bandi2023} and Nti et al.~\cite{nti2023} each reported
classification accuracies above 99\% using voting and tree-based ensembles
respectively -- a pattern of very high reported accuracy across this literature that
underscores why uncertainty and generalizability, not just point accuracy, matter for
practical deployment. Gopi and Karthikeyan~\cite{gopi2024} combined Red Fox
Optimization with an ensemble recurrent neural network for joint crop recommendation
and yield prediction. Elbasi et al.~\cite{elbasi2023} compared fifteen ML algorithms
combined with IoT sensor data, reporting 97\% accuracy but, like most prior work, only
deterministic outputs. Bera et al.~\cite{bera2023} proposed an edge/dew-computing and
federated-learning architecture (E-CropReco) to reduce latency in real-time crop
recommendation. Thilakarathne et al.~\cite{thilakarathne2023} built a cloud-based IoT
platform for real-time monitoring, though scoped to a single crop (tomato). Shingade
et al.~\cite{shingade2022} developed a Random Forest classifier specifically for seed
recommendation. Senapaty et al.~\cite{senapaty2024} and Chaudhari et
al.~\cite{chaudhari2025} each describe decision-support-oriented crop recommendation
systems closer in spirit to a deployed tool than an offline benchmark, and Agrawal et
al.~\cite{agrawal2024} combined ML with spatial multi-criteria decision analysis for
agricultural land suitability classification, one of the few studies in this space
that, like the present work, reasons at the level of a specific geographic unit
rather than a generic feature vector. Afzal et al.~\cite{afzal2025} combined Random
Forest and XGBoost with soil information to improve recommendation output.

A separate line of work analyzes rainfall trends directly rather than using rainfall
only as a prediction feature. Praveen et al.~\cite{praveen2020} applied the
Mann-Kendall trend test, Pettitt change-point detection, and an ANN-MLP forecaster to
115 years (1901--2015) of Indian meteorological-division rainfall data, finding
statistically significant negative rainfall trends in most divisions -- direct
evidence that Indian monsoon behavior is measurably shifting over the long term, which
motivates conditioning a crop recommendation on a district's own historical trend
rather than its current-season rainfall alone, as this paper does (Section IV).
Correndo et al.~\cite{correndo2021} combined Bayesian statistics and machine learning
to show that weather is the dominant source of uncertainty in maize yield response to
nitrogen, reinforcing that climate variability -- not just soil or management inputs
-- is a primary driver of yield uncertainty. Rihan et al.~\cite{rihan2023} used
machine learning and deep learning with explicit sensitivity and uncertainty analysis
for forest-fire susceptibility mapping and identified rainfall as one of the dominant
driving factors, illustrating that rainfall-trend-aware risk modeling is valuable
beyond crop recommendation specifically.

Beyond crop \emph{recommendation}, a closely related line of work addresses yield
\emph{prediction} directly with tree-ensemble methods. Jeong et al.~\cite{jeong2016}
evaluated Random Forests against multiple linear regression for wheat, maize, and
potato yield prediction at global and regional scale, reporting RMSE of 6--14\% of
observed yield for Random Forest versus 14--49\% for the linear baseline -- an early,
influential demonstration that tree ensembles capture climate-yield relationships that
linear models miss, consistent with the gap we observe between our single Decision
Tree baseline and the ensemble methods in Section V. Hu et al.~\cite{hu2023}
specifically cautioned against black-box ML models for evaluating climate-change
impacts on crop yield and argued for explainable, interpretable approaches instead --
a concern this paper addresses in part through per-feature importance analysis
(Section V-B) rather than treating the trained model as opaque.

These studies share two characteristics relevant to the present work. First, nearly
all are evaluated on generic or single-crop datasets rather than a real administrative
region's multi-decade agricultural census, so a trained model has no direct notion of
\emph{which named district, which government-defined season, and which historical
climate trend} apply to a given farmer. Second, with the exception of
\cite{alam2025}, none quantify how confident a recommendation is, and even
\cite{alam2025} does not carry its model through to a farmer-facing deployed
application. The system described in this paper addresses the first gap directly by
building on real district-level government data for Tamil Nadu and deploying the
result as an installable web and Android application; we discuss the second gap --
uncertainty quantification -- as future work in Section VII, for which the
entropy-based approach of \cite{alam2025} is a natural candidate to adapt to a
regression setting.

\section{Data}

\subsection{Source and Scope}

The system is built on two merged real data sources: (i) a government agricultural
production census covering Tamil Nadu, reporting district, year, season, crop, area
under cultivation, and production, for the years 2005--2019; and (ii) district-level
historical climate records (temperature, rainfall, humidity, wind) for the same
period and districts. The two sources were merged on district, year, and season to
produce a single feature table used for both model training and the application's
real-time climate-retrieval logic.

The application supports all 38 present-day Tamil Nadu districts. Five districts
created after the 2019 administrative reorganization (Kallakurichi, Mayiladuthurai,
Ranipet, Tenkasi, Tirupathur) have no independent history under their current name in
the census data; these are mapped to the parent district they were carved from for
the purposes of historical lookups, while remaining independently selectable in the
application interface.

\subsection{Crop Selection Criterion}

The government census records production for many more crops than the eight
originally covered by an early version of this system (Rice, Groundnut, Cotton,
Sugarcane, Maize, Bajra, Ragi, Banana). To responsibly expand coverage, every
additional crop reported for Tamil Nadu in the source census was evaluated against a
fixed data-sufficiency bar and included only if it met \emph{all three} of the
following conditions, chosen to match the same order of magnitude as the original
eight crops:
\begin{itemize}
  \item at least 300 recorded (district, year, season) instances;
  \item at least 15 distinct years of history;
  \item at least 15 distinct districts represented.
\end{itemize}

Seventeen crops cleared this bar and were added: five pulses (Urad, Moong,
Horse-gram, Arhar/Tur, Gram), two oilseeds (Sesamum, Sunflower), two millets (Jowar,
Small millets), three spices (Turmeric, Dry chillies, Coriander), one plantation crop
(Cashewnut), and four others (Tobacco, Onion, Tapioca, Sweet potato), bringing total
coverage to 25 crops. Crops that failed the bar were deliberately excluded rather than
included with unreliable history: Mango, Grapes, Papaya, and Citrus had only two years
of recorded data each; Korra, Varagu, Samai, Mesta, and Sannhamp had two to three
years each; Potato, Garlic, Ginger, Black pepper, and Cardamom fell under the
record/district threshold. Coconut was excluded for a different reason described in
the next subsection. Tea, Coffee, and Rubber are not present in this census source at
all -- they are tracked by separate commodity boards in India -- and would require an
entirely different data source to add responsibly.

\subsection{A Systematic Production-Unit Inconsistency}

While computing per-crop average yield for the application's crop-information
feature, yield figures for individual (district, crop, year) combinations were found
to be physically implausible for a subset of records -- e.g., rice in one district
reporting approximately 3--4 tonnes/hectare for 2005--2013 and then approximately
330--490 tonnes/hectare for 2014--2019, an impossible jump for the same crop and
district. Manual inspection of the raw census source revealed that many
(district, crop, year, season) combinations are recorded \emph{twice}, with one copy's
Production value exactly $100\times$ the other (confirmed on thousands of matching
duplicate-row pairs, with ratios consistently in the 90--110$\times$ range). An
existing deduplication step in the data pipeline already preferred the smaller of two
duplicate values when both copies of a record existed, but a large share of records
-- especially from 2014 onward, and for some districts in every year -- had only the
inflated copy submitted to the source, leaving nothing for the deduplication step to
prefer.

We verified that this bimodal pattern is present in every one of the 25 crops, not
only rice, by examining the quantile distribution of the raw Yield column per crop:
every crop showed a low cluster consistent with plausible tonnes/hectare values and a
high cluster approximately $100\times$ larger. A purely statistical correction (e.g.,
the largest gap in the sorted per-crop value distribution) was evaluated and rejected,
because it let some crops' genuinely implausible values through as ``legitimate'' --
for instance, it would have accepted Groundnut values up to 14.93 t/ha, far beyond the
world-record groundnut yield of approximately 4--5 t/ha.

Instead, we derived a correction empirically from the confirmed duplicate pairs
themselves. For each crop, we collected every confirmed duplicate pair (production
ratio in the 90--110$\times$ range) and computed the 99th percentile of the
\emph{correct}-scale member of each pair. Each crop's correction ceiling was set to
$1.5\times$ this percentile -- comfortably above every observed genuine value and
comfortably below the $\sim$100$\times$-inflated ones. Any record whose implied yield
(Production $\div$ Area) exceeded its crop's ceiling had its Production value divided
by 100. As a safety net, any record still more than $3\times$ its ceiling after this
correction was dropped rather than further guessed at; this caught exactly one record
(a Cashewnut observation with an Area value of 1.0 hectare, an unrelated broken input
that no unit correction can fix). In total, 6{,}295 records were rescaled and one
record was dropped across the full 25-crop dataset. Separately, Coconut was found to
have the same duplicate-record pattern but with ratios near $100{,}000\times$ rather
than a clean $\sim$100$\times$, indicating a different and unverified error mode; it
was excluded from the crop list rather than corrected without equivalent evidence.

We additionally found that the per-crop yield-tertile categorization used elsewhere in
the pipeline (a Low/Medium/High label derived from global yield quantiles) exhibited a
``crop-identity leak'': because raw yield scales differ enormously by crop (e.g.,
Sesamum's typical yield is on the order of 0.5 t/ha versus Sugarcane's $\sim$100 t/ha),
a single global quantile split classified nearly every Sugarcane record as ``High''
and nearly every Sesamum record as ``Low'' regardless of how each performed relative
to its own crop's normal range. This was corrected by computing yield quantiles
\emph{per crop} rather than globally, after which every crop's category distribution
split into roughly even thirds.

\section{Monsoon Shift Analysis and Recommendation Pipeline}

Given a district, crop, and season, the recommendation pipeline executes five stages:

\begin{enumerate}
\item \textbf{Climate data retrieval}: historical climate and agriculture features
      for the exact (district, crop, season) combination are averaged from the
      corrected dataset, falling back to (district, season), then district-only, then
      (crop, season), then the full dataset, if the exact combination has insufficient
      history.
\item \textbf{Yield prediction}: the trained regression model (Section V) predicts
      yield in tonnes/hectare from twenty features -- district, year, season, crop,
      cultivated area, six raw climate statistics (average/max/min temperature, total
      and average rainfall, average wind speed), three derived statistics (rainy days,
      temperature and rainfall standard deviation), and four categorical bucket
      features (temperature, rainfall, humidity, and wind category) -- and the
      predicted value is bucketed into a per-crop yield tertile (Section III-C).
\item \textbf{Monsoon-shift-aware climate risk assessment}: drought and flood risk are
      derived jointly from the district's current-season rainfall category and its
      real 21-year historical rainfall trend (percentage change and direction),
      rather than from the current season's rainfall alone. This mirrors the
      trend-then-risk pattern used in Indian rainfall trend-detection studies such as
      \cite{praveen2020}, which found significant long-term negative rainfall trends
      across most Indian meteorological divisions, and in rainfall-driven hazard
      studies such as \cite{rihan2023}, which identified rainfall as a dominant driver
      of susceptibility alongside evapotranspiration; here the trend is computed
      per-district from the corrected historical series (Section III-C) rather than
      at the coarser meteorological-division scale, and is used to shift the
      drought/flood classification a district would otherwise receive from its
      current season's rainfall category alone -- e.g., a district reporting
      currently-adequate rainfall but a strong long-term declining trend is flagged at
      elevated drought risk rather than ``Low.''
\item \textbf{Sowing window selection}: a TNAU-style (Tamil Nadu Agricultural
      University) reference calendar maps each (crop, season) pair to an agronomic
      sowing-window date range, built from established Tamil Nadu agronomic-extension
      practice. Because the underlying dataset contains no sowing-date ground truth,
      this stage relies on domain-expert reference values rather than a learned date.
      If the current year's window for the requested crop and season has already
      elapsed, the system correctly advances the recommendation to the next
      occurrence rather than recommending an already-passed date, and surfaces this
      explicitly to the farmer.
\item \textbf{Confidence scoring}: a data-driven confidence score is derived from the
      predicted yield's tertile category and the overall drought/flood risk level,
      rather than a fixed constant.
\end{enumerate}

\section{Model Comparison}

\subsection{Setup}

Three regression models were trained on the corrected 25-crop feature table with an
80:20 train/test split (\texttt{random\_state=42}): a single Decision Tree regressor
as a baseline, a bagged Random Forest regressor~\cite{biau2016}
(\texttt{n\_estimators=300}, \texttt{max\_depth=20}, \texttt{min\_samples\_split=5},
\texttt{min\_samples\_leaf=2}), and a gradient-boosted XGBoost regressor~\cite{chen2016}.

An initial XGBoost configuration used substantially more conservative hyperparameters
than the Random Forest baseline (\texttt{max\_depth=3}, \texttt{n\_estimators=300},
strong L1/L2 regularization), which is not a fair comparison: Random Forest was given
free rein to depth 20 while XGBoost was deliberately constrained to depth 3. Under
this mismatched configuration, XGBoost under-performed Random Forest. We therefore
re-tuned XGBoost to a comparable trees $\times$ depth budget
(\texttt{n\_estimators=800}, \texttt{max\_depth=6}, \texttt{learning\_rate=0.03},
\texttt{min\_child\_weight=1}, \texttt{subsample=0.85}, \texttt{colsample\_bytree=0.85},
\texttt{reg\_alpha=0.01}, \texttt{reg\_lambda=1.0}), selected by checking both
held-out test $R^2$ and 5-fold cross-validated $R^2$ together, rather than picking
whichever configuration scored highest on a single lucky split.

\subsection{Results}

Table~\ref{tab:results} reports held-out test performance for all three models after
the fair re-tuning. XGBoost achieves the best test $R^2$ (96.98\%), ahead of Random
Forest (96.57\%) and the single Decision Tree (95.27\%), with the lowest MAE and RMSE
among the three. Fig.~\ref{fig:comparison} shows the same comparison graphically.

\begin{table}[htbp]
\caption{Model Comparison on the Corrected Dataset (Held-Out Test Split)}
\label{tab:results}
\centering
\begin{tabular}{lccccc}
\toprule
Model & Train $R^2$ & Test $R^2$ & MAE & RMSE & Gap \\
\midrule
XGBoost & 99.60\% & \textbf{96.98\%} & \textbf{1.45} & \textbf{3.82} & 0.026 \\
Random Forest & 98.38\% & 96.57\% & 1.19 & 4.08 & 0.018 \\
Decision Tree & 97.60\% & 95.27\% & 1.46 & 4.79 & 0.023 \\
\bottomrule
\end{tabular}
\end{table}

\begin{figure}[htbp]
\centering
\includegraphics[width=0.95\linewidth]{figures/model_comparison_bar_chart.png}
\caption{Test $R^2$ comparison across the three regression models.}
\label{fig:comparison}
\end{figure}

Random Forest has the lowest overfitting gap (train $R^2$ minus test $R^2$) of the
three models (0.018 versus XGBoost's 0.026), which raises a natural question: does a
smaller gap make it the better model? We argue no, for two reasons. First, both gaps
are comfortably inside what we treat as a ``good generalization'' threshold (gap
$\leq 0.05$); neither model shows concerning overfitting. Second, overfitting gap
measures the \emph{consistency} of a model's score between training and test data, not
its absolute quality -- a model that always predicted the per-crop mean yield would
have a gap near zero while being useless. Once both models are confirmed to
generalize well, the appropriate tie-breaker is which one is more accurate on unseen
data, which is what test $R^2$ measures directly, and that favors XGBoost. XGBoost's
slightly larger gap is also structurally expected: boosting fits training data more
closely by design (each tree corrects the residual of the previous ones), whereas
bagging's averaging across independently-trained trees smooths training fit further --
a difference in mechanism, not a flaw in either model.

To verify this result was not an artifact of one lucky 80:20 split, we additionally
ran 5-fold cross-validation: the tuned XGBoost configuration achieved a mean CV
$R^2$ of 0.9585 versus Random Forest's 0.9576 and the original, under-tuned XGBoost
configuration's 0.9467 -- confirming that fair tuning, not the particular split,
drives the improvement, though the cross-validated margin between XGBoost and Random
Forest (about 0.1 percentage points) is markedly smaller than the held-out single-split
margin (about 0.4 points), indicating the true difference between the two ensemble
methods on this dataset is modest.

Fig.~\ref{fig:importance} shows XGBoost's feature importance. Consistent with
agronomic expectation, climate variables (rainfall and temperature statistics) and
cultivated area dominate, while the categorical bucket features contribute
comparatively little additional signal once their underlying continuous values are
already present.

\begin{figure}[htbp]
\centering
\includegraphics[width=0.95\linewidth]{figures/xgboost_feature_importance_bar.png}
\caption{XGBoost feature importance for yield prediction.}
\label{fig:importance}
\end{figure}

\section{Deployed System}

The trained model is served through a full-stack web application rather than remaining
an offline notebook artifact. The backend (Python, Flask) exposes the five-stage
pipeline of Section IV as a farmer-facing dashboard, a district/crop/season
recommendation form, a monsoon/climate-trend view, a sowing calendar view, and a
notification center, in both English and Tamil. Application data (farmer accounts and
stored recommendations) is persisted in PostgreSQL. The dashboard's ``today's
weather'' figures are retrieved live per district from the Open-Meteo API at request
time (cached for fifteen minutes), replacing an earlier static per-district estimate
that farmers correctly identified as inconsistent with real observed conditions.

The same application is exposed as an installable Progressive Web App (a web manifest
and service worker enable ``Add to Home Screen'' installation with offline caching of
static assets) and, separately, as a native Android application built with Capacitor,
wrapping the same served pages in a WebView shell rather than re-implementing the
interface in a native toolkit -- preserving a single source of truth for the
interface across the browser, PWA, and Android targets. The system is deployed to a
cloud host with a managed PostgreSQL instance, so the application is reachable from
any network, not only a local development machine.

\section{Discussion, Limitations, and Future Work}

The production-unit inconsistency described in Section III-C is, to our knowledge, not
previously documented for this class of Indian government agricultural census data,
and it has a direct implication for any prior or future model trained on similar
uncorrected sources: a substantial fraction of the yield target variable may be
corrupted by a consistent, detectable, but easily-overlooked unit error. We consider
the detection and correction procedure -- grounding the correction in confirmed
duplicate-record evidence rather than a purely statistical heuristic -- to be as
significant a contribution of this work as the model comparison itself, since it
determines whether the regression target is meaningful at all.

Several limitations remain. First, following the discussion in Section II, this
system does not currently quantify the confidence of a given yield prediction or
sowing recommendation beyond the coarse per-crop tertile category and a heuristic
confidence score; adapting an entropy-based or ensemble-variance-based uncertainty
measure, in the style of \cite{alam2025}, to this regression setting is a natural next
step, particularly for districts or crops with sparse historical records where the
model's fallback averaging (Section IV, stage 1) already signals reduced data
support. Second, the historical data underlying both the yield model and the climate
risk assessment ends in 2019; incorporating more recent government releases as they
become available, and integrating real-time soil-sensor data where farmers have
access to it, would improve currency. Third, the sowing-window calendar is a
domain-expert reference rather than a learned quantity, since the underlying dataset
contains no sowing-date ground truth; a future direction is to source sowing-date
records (e.g., from agricultural extension offices) and let the model learn
district-specific timing shifts directly. Finally, this system is scoped to Tamil
Nadu; extending the same data-correction and modeling methodology to other Indian
states' government agricultural census data is a direct avenue for broader impact,
echoing the generalization concern raised for the generic benchmark dataset in
Section II.

\section{Conclusion}

This paper presented a district-level crop sowing recommendation system built on real
Tamil Nadu government agricultural and climate records, expanded from an initial eight
crops to twenty-five under an explicit, defensible data-sufficiency criterion. In the
course of building it, we identified and corrected a systematic $\sim$100$\times$
production-unit inconsistency present in the underlying government source for every
crop -- a finding with direct implications for the validity of yield models trained on
this and similar data sources -- and corrected a related per-crop yield-categorization
flaw. On the corrected data, a fairly-tuned XGBoost regressor achieved the best
held-out performance (test $R^2 = 96.98\%$) among Decision Tree, Random Forest, and
XGBoost baselines, verified with cross-validation. The resulting five-stage
recommendation pipeline is deployed as a full-stack, multilingual web application with
live weather integration, installable as both a Progressive Web App and a native
Android application, moving beyond the offline, single-dataset evaluation common in
prior crop recommendation literature toward a system a farmer can actually open on
their phone today. Uncertainty quantification, more recent and sowing-date-labeled
data, and generalization to other states remain the primary directions for future
work.

\begin{thebibliography}{00}

\bibitem{alam2025} Md. S. B. Alam, V. Esichaikul, A. Lameesa, S. F. Ahmed, and
A. H. Gandomi, ``An approach for crop recommendation with uncertainty quantification
based on machine learning for sustainable agricultural decision-making,'' \emph{Results
in Engineering}, vol. 26, p. 105505, 2025.

\bibitem{elbasi2023} E. Elbasi, C. Zaki, A. E. Topcu, W. Abdelbaki, A. I. Zreikat,
E. Cina, A. Shdefat, and L. Saker, ``Crop prediction model using machine learning
algorithms,'' \emph{Applied Sciences}, vol. 13, no. 16, p. 9288, 2023.

\bibitem{bera2023} S. Bera, T. Dey, A. Mukherjee, and R. Buyya, ``E-CropReco: a
dew-edge-based multi-parametric crop recommendation framework for internet of
agricultural things,'' \emph{The Journal of Supercomputing}, vol. 79, 2023.

\bibitem{thilakarathne2023} N. N. Thilakarathne, M. S. A. Bakar, P. E. Abas, and
H. Yassin, ``Towards making the fields talk: a real-time cloud enabled IoT crop
management platform for smart agriculture,'' \emph{Frontiers in Plant Science},
vol. 13, 2023.

\bibitem{afzal2025} H. Afzal, M. Amjad, A. Raza, K. Munir, S. G. Villar,
L. A. D. Lopez, and I. Ashraf, ``Incorporating soil information with machine learning
for crop recommendation to improve agricultural output,'' \emph{Scientific Reports},
vol. 15, 2025.

\bibitem{chen2016} T. Chen and C. Guestrin, ``XGBoost: a scalable tree boosting
system,'' in \emph{Proc. ACM SIGKDD Int. Conf. Knowledge Discovery and Data Mining},
2016.

\bibitem{biau2016} G. Biau and E. Scornet, ``A random forest guided tour,''
\emph{Test}, vol. 25, 2016.

\bibitem{jeong2016} J. H. Jeong, J. P. Resop, N. D. Mueller, D. H. Fleisher, K. Yun,
E. E. Butler, D. J. Timlin, K.-M. Shim, J. S. Gerber, V. R. Reddy, and S.-H. Kim,
``Random forests for global and regional crop yield predictions,'' \emph{PLOS ONE},
vol. 11, no. 6, p. e0156571, 2016.

\bibitem{hu2023} T. Hu, X. Zhang, G. Bohrer, Y. Liu, Y. Zhou, J. Martin, Y. Li, and
K. Zhao, ``Crop yield prediction via explainable AI and interpretable machine
learning: dangers of black box models for evaluating climate change impacts on crop
yield,'' \emph{Agricultural and Forest Meteorology}, vol. 336, p. 109458, 2023.

\bibitem{praveen2020} B. Praveen, S. Talukdar, Shahfahad, S. Mahato, J. Mondal,
P. Sharma, A. R. M. T. Islam, and A. Rahman, ``Analyzing trend and forecasting of
rainfall changes in India using non-parametrical and machine learning approaches,''
\emph{Scientific Reports}, vol. 10, p. 10342, 2020.

\bibitem{hasan2023} M. Hasan, M. A. Marjan, M. P. Uddin, M. I. Afjal, S. Kardy, S. Ma,
and Y. Nam, ``Ensemble machine learning-based recommendation system for effective
prediction of suitable agricultural crop cultivation,'' \emph{Frontiers in Plant
Science}, vol. 14, 2023.

\bibitem{bandi2023} R. Bandi, M. S. S. Likhit, S. R. Reddy, S. R. Bodla, and
V. S. Venkat, ``Voting classifier-based crop recommendation,'' \emph{SN Computer
Science}, vol. 4, 2023.

\bibitem{nti2023} I. K. Nti, A. Zaman, O. Nyarko-Boateng, A. F. Adekoya, and
F. Keyeremeh, ``A predictive analytics model for crop suitability and productivity
with tree-based ensemble learning,'' \emph{Decision Analytics Journal}, vol. 8, 2023.

\bibitem{gopi2024} P. S. S. Gopi and M. Karthikeyan, ``Red fox optimization with
ensemble recurrent neural network for crop recommendation and yield prediction
model,'' \emph{Multimedia Tools and Applications}, vol. 83, 2024.

\bibitem{shingade2022} S. D. Shingade, R. P. Mudhalwadkar, and K. M. Masal, ``Random
forest machine learning classifier for seed recommendation,'' in \emph{Proc. Int.
Conf. Edge Computing and Applications (ICECAA)}, 2022.

\bibitem{senapaty2024} M. K. Senapaty, A. Ray, and N. Padhy, ``A decision support
system for crop recommendation using machine learning classification algorithms,''
\emph{Agriculture}, vol. 14, no. 8, p. 1256, 2024.

\bibitem{chaudhari2025} R. R. Chaudhari, S. Jain, and S. Gupta, ``Agricultural
machine learning platform: enhancing crop suggestion and crop yield estimates,''
\emph{Journal of Integrated Science and Technology}, vol. 13, 2025.

\bibitem{agrawal2024} N. Agrawal, H. Govil, and T. Kumar, ``Agricultural land
suitability classification and crop suggestion using machine learning and spatial
multicriteria decision analysis in semi-arid ecosystem,'' \emph{Environment,
Development and Sustainability}, 2024.

\bibitem{correndo2021} A. A. Correndo, N. Tremblay, J. A. Coulter, D. Ruiz-Diaz,
D. Franzen, E. Nafziger, V. Prasad, L. H. M. Rosso, K. Steinke, J. Du, C. D. Messina,
and I. A. Ciampitti, ``Unraveling uncertainty drivers of the maize yield response to
nitrogen: a Bayesian and machine learning approach,'' \emph{Agricultural and Forest
Meteorology}, vol. 311, p. 108668, 2021.

\bibitem{rihan2023} M. Rihan, A. Ali Bindajam, S. Talukdar, Shahfahad,
M. W. Naikoo, J. Mallick, and A. Rahman, ``Forest fire susceptibility mapping with
sensitivity and uncertainty analysis using machine learning and deep learning
algorithms,'' \emph{Advances in Space Research}, vol. 72, 2023.

\end{thebibliography}

\end{document}
